%==============================================================================
% Capitulo 1: Introduccion
%==============================================================================
\section{Introduccion}

\subsection{Contextualizacion}

La evolucion de las tecnologias de la informacion ha transformado profundamente la dinamica operativa de las organizaciones contemporaneas. En la actualidad los sistemas informaticos constituyen el soporte fundamental para la ejecucion de procesos administrativos, productivos y de comunicacion, y esa dependencia tecnologica ha incrementado de manera correlativa la importancia de los servicios de soporte tecnico, particularmente aquellos relacionados con la gestion de incidentes informaticos. \textcite{Galup2009_itsm} caracterizan a la disciplina de gestion de servicios de tecnologias de la informacion (ITSM, por su sigla en ingles) como un cuerpo de practicas orientado a alinear la operacion tecnologica con los objetivos del negocio, donde la mesa de ayuda actua como punto unico de contacto entre el usuario y el area de tecnologia.

Dentro de este escenario, las mesas de ayuda empresariales cumplen un rol central en la recepcion, registro y seguimiento de incidentes reportados por los usuarios. Estas unidades organizacionales canalizan solicitudes y constituyen la primera linea de atencion tecnica, de modo que su eficiencia impacta directamente en la continuidad operativa. No obstante, en muchas organizaciones medianas el proceso de registro inicial continua realizandose manualmente, lo que introduce demoras, inconsistencias y una marcada dependencia del criterio individual del operador. \textcite{Galup2009_itsm} advierten que los procesos manuales de soporte tienden a degradarse cuando el volumen de solicitudes crece de manera no lineal o cuando el personal debe alternar entre tareas operativas heterogeneas.

El contexto argentino presenta caracteristicas particulares que acentuan esta problematica. Las organizaciones medianas del pais, definidas por la Secretaria de la Pequena y Mediana Empresa como aquellas que emplean entre 50 y 200 trabajadores en el sector servicios, operan con presupuestos de tecnologia acotados y enfrentan restricciones cambiarias que encarecen la adopcion de soluciones comerciales internacionales. La combinacion de estas restricciones con una demanda creciente de servicios digitales configura un escenario en el que las soluciones basadas en componentes de codigo abierto y servicios de inferencia bajo demanda resultan particularmente atractivas desde el punto de vista economico y operativo.

\subsection{Planteamiento del problema}

La presente investigacion se desarrolla en el marco de una organizacion mediana del sector servicios con sede en la provincia de Mendoza, Republica Argentina, cuya plantilla asciende a aproximadamente 120 usuarios internos distribuidos en cinco areas funcionales. El sector de Operaciones desempena simultaneamente tareas vinculadas a la gestion operativa cotidiana y a la atencion de incidentes informaticos reportados por los usuarios, lo que obliga al personal responsable a interrumpir actividades criticas para registrar manualmente fallas de hardware, inconvenientes de software o solicitudes de asistencia tecnica.

Las mediciones preliminares realizadas durante una semana laboral representativa evidenciaron un volumen promedio de 42 incidentes diarios, un tiempo medio de registro manual de 2 minutos con 45 segundos por incidente y una tasa de derivacion erronea inicial cercana al 15~\%. La proyeccion trimestral sobre estos indicadores arrojo un universo aproximado de 3.700 incidentes, con un costo operativo asociado a la etapa de registro estimado en el orden de 200 horas-persona por trimestre. Este tiempo, que el personal podria redirigir a tareas de mayor valor agregado si la etapa de registro estuviera automatizada, representa un costo de oportunidad significativo para la organizacion.

El registro manual de incidentes implica ademas la lectura de correos electronicos, la interpretacion de descripciones textuales heterogeneas, la clasificacion del problema segun el area responsable y la carga manual de la informacion en un sistema de tickets. Este procedimiento requiere tiempo, interrumpe las tareas del personal tecnico y puede generar errores de clasificacion que prolongan la cadena de resolucion cuando un ticket es derivado a un area incorrecta. La dispersion de los canales de comunicacion ---correo, llamada telefonica y formularios internos--- agrava la situacion al obligar al personal a monitorear multiples fuentes simultaneamente.

La naturaleza del problema presenta un patron recurrente en la literatura de gestion de servicios: un cuello de botella en la etapa inicial del proceso que degrada la eficiencia del ciclo completo. \textcite{Galup2009_itsm} documentan que los procesos de soporte no automatizados tienden a degradarse cuando el volumen de solicitudes crece, con un impacto desproporcionado en organizaciones medianas donde el personal de soporte combina multiples roles. Estos indicadores justifican la busqueda de una alternativa tecnologica que reduzca el costo operativo, mejore la calidad de la derivacion temprana y disminuya la variabilidad asociada al criterio individual.

\subsection{Pregunta de investigacion}

A partir del problema descrito, esta investigacion se orienta a responder la siguiente pregunta: ?`reduce la automatizacion del registro inicial de incidentes mediante orquestacion de flujos y procesamiento de lenguaje natural el tiempo de registro y mejora la exactitud de clasificacion en una mesa de ayuda empresarial mediana, en comparacion con el proceso manual realizado por personal experimentado?

La pregunta se descompone en tres interrogantes subsidiarios que guian el diseno experimental y la seleccion de metricas. Primero, ?`cual es la magnitud de la reduccion del tiempo de registro que puede obtenerse mediante la automatizacion, y resulta esa reduccion estadisticamente significativa? Segundo, ?`alcanza el clasificador automatico una exactitud comparable o superior a la del criterio humano, y como se distribuyen los errores entre las categorias objetivo? Tercero, ?`permite la arquitectura propuesta integrar canales multiples de recepcion sin degradar la latencia ni la calidad de la clasificacion?

\subsection{Hipotesis de trabajo}

La hipotesis principal de este trabajo sostiene que la implementacion de un sistema automatizado de registro de incidentes basado en orquestacion de flujos con N8N y procesamiento de lenguaje natural mediante un modelo de lenguaje grande permite reducir significativamente el tiempo de registro y la proporcion de intervencion humana, manteniendo una exactitud de clasificacion equivalente o superior a la observada en el proceso manual realizado por personal experimentado.

Formalmente, la hipotesis nula \(H_{0}\) postula la inexistencia de diferencias significativas entre ambos flujos en las variables observadas: \(H_{0}: M_{\text{manual}} = M_{\text{automatizado}}\), donde \(M\) representa la mediana poblacional del tiempo de registro. La hipotesis alternativa \(H_{1}\) postula que la mediana del flujo automatizado es inferior: \(H_{1}: M_{\text{manual}} > M_{\text{automatizado}}\).

Como hipotesis subsidiarias se plantea, en primer lugar, que la integracion de multiples canales de entrada dentro de un unico flujo automatizado mejora la accesibilidad del sistema y reduce la dispersion operativa sin introducir cuellos de botella adicionales. En segundo lugar, que un esquema hibrido de clasificacion basado en reglas deterministicas y modelo de lenguaje grande resulta mas eficiente ---en terminos de la relacion entre exactitud, latencia y costo de inferencia--- que un esquema construido exclusivamente sobre modelo externo, dado que una proporcion sustancial de los incidentes presenta marcadores lexicos inequivocos que no requieren consulta al modelo de lenguaje.

\subsection{Objetivo general}

Desarrollar e implementar un sistema automatizado de mesa de ayuda empresarial basado en orquestacion de flujos con N8N y procesamiento de lenguaje natural en Python que permita recibir, procesar, clasificar y registrar incidentes de usuarios de manera automatica, reduciendo el tiempo de registro y la proporcion de intervencion humana en la etapa inicial del proceso, manteniendo la supervision tecnica para la resolucion de los casos y garantizando la trazabilidad de las decisiones tomadas por el sistema.

\subsection{Objetivos especificos}

El primer objetivo especifico consiste en disenar una arquitectura distribuida y modular que integre un motor de orquestacion de flujos, un modulo de procesamiento de lenguaje natural, un canal de telefonia con transcripcion automatica del habla y una base de datos relacional, garantizando comunicacion asincronica entre componentes mediante interfaces de programacion de aplicaciones de tipo REST sobre HTTP cifrado.

El segundo objetivo especifico se orienta a implementar un clasificador automatico que asigne cada incidente a uno de los tres sectores responsables ---Sistemas, Operaciones o Soporte Tecnico--- alcanzando una exactitud global no inferior al 85~\% sobre un corpus controlado de 200 casos representativos del entorno operativo, y un valor F1 macro promediado no inferior a 0,85. El umbral del 85~\% se fundamenta en tres criterios convergentes. Primero, los trabajos de \textcite{Paramesh2019_it_service_desk} reportan exactitudes de aproximadamente 87~\% con enfoques de maquinas de vectores de soporte (SVM) y representaciones TF-IDF, estableciendo un piso de referencia para el estado del arte en clasificacion de tickets. Segundo, el costo organizacional de una derivacion erronea implica al menos 10 a 15 minutos adicionales de reasignacion segun las mediciones preliminares de la organizacion, por lo que una exactitud inferior al 85~\% resultaria en tiempos de resolucion comparables a los del flujo manual para un volumen de errores inaceptable. Tercero, el margen por encima del 85~\% reserva espacio estadistico para que el intervalo de confianza del estimador no solape la zona de rechazo de la hipotesis, dado el tamano muestral de 200 casos.

El tercer objetivo especifico plantea integrar tres canales de entrada paralelos ---correo electronico, formulario web y llamada telefonica con transcripcion automatica--- dentro de un unico flujo automatizado que normalice la entrada y centralice el procesamiento, mejorando la accesibilidad del sistema sin incrementar la complejidad operativa percibida por los usuarios finales.

El cuarto objetivo especifico consiste en evaluar comparativamente el flujo manual y el flujo automatizado en terminos de tiempo medio de registro, dispersion, exactitud de clasificacion y proporcion de intervencion humana, mediante un diseno cuasiexperimental con muestras pareadas, prueba estadistica no parametrica y reporte de intervalos de confianza al 95~\%. El diseno experimental sigue las recomendaciones de \textcite{HernandezSampieri2018_metodologia} para investigaciones tecnologicas de caracter aplicado.

El quinto objetivo especifico se orienta a documentar exhaustivamente la arquitectura, los procedimientos de despliegue y los resultados experimentales, proponiendo una hoja de ruta de evolucion que oriente futuras mejoras y replicaciones del estudio en organizaciones de caracteristicas comparables.

\subsection{Justificacion}

La justificacion del trabajo se sostiene en tres dimensiones complementarias.

Desde la dimension organizacional, la automatizacion propuesta libera capacidad operativa del personal tecnico, reduce el tiempo de respuesta percibido por los usuarios y disminuye la variabilidad introducida por el criterio individual. La evidencia recogida en este trabajo permite cuantificar esa mejora en el contexto especifico de una organizacion mediana argentina, aportando un caso documentado que otras organizaciones pueden utilizar como referencia para estimar el retorno esperado de una inversion similar.

Desde la dimension tecnologica, la solucion integra herramientas de bajo costo y codigo abierto ---N8N, Python, FastAPI, PostgreSQL--- permitiendo una implementacion reproducible en organizaciones de tamano mediano sin requerir grandes inversiones de capital ni dependencia exclusiva de proveedores comerciales. Esta propiedad es particularmente relevante en el contexto argentino, donde las restricciones de acceso a divisas y los costos de suscripcion en moneda extranjera constituyen barreras significativas para la adopcion de soluciones comerciales internacionales. Ademas, el despliegue autoalojado de la totalidad de los componentes ---con excepcion del servicio de inferencia--- otorga a la organizacion control pleno sobre sus datos operativos, en linea con el principio de soberania digital promovido por la Ley~25.326 de Proteccion de los Datos Personales \parencite{Ley25326_2000}.

Desde la dimension academica, el trabajo aporta evidencia empirica sobre el desempeno de modelos de lenguaje grandes aplicados a la clasificacion de tickets en espanol rioplatense, dominio escasamente documentado en la literatura comparativa disponible. La mayoria de los estudios publicados sobre clasificacion automatica de tickets utilizan corpus en ingles, aleman o portugues, y los pocos trabajos que incluyen el espanol lo hacen con variantes peninsulares o con corpus genericos no especificos del dominio de mesas de ayuda. Este trabajo contribuye a cerrar esa brecha mediante la validacion rigurosa de un clasificador hibrido sobre un corpus de 200 casos en espanol rioplatense con doble etiquetado independiente y acuerdo inter-evaluador cuantificado.

Adicionalmente, la mayoria de las soluciones comerciales disponibles ---tales como Zendesk Suite, Freshdesk o ServiceNow ITSM--- presentan modelos de costos por agente y mes que escalan linealmente con el equipo y operan exclusivamente en modalidad de software como servicio con almacenamiento de datos en jurisdicciones extranjeras. La propuesta desarrollada en esta investigacion utiliza herramientas de despliegue autoalojado, lo que facilita su adopcion en organizaciones medianas y simplifica el cumplimiento de la normativa argentina de proteccion de datos personales.

\subsection{Alcance y delimitaciones}

El alcance del trabajo se circunscribe a la etapa de recepcion, clasificacion y registro inicial de incidentes, sin abarcar la resolucion tecnica de los mismos, la cual permanece bajo responsabilidad del personal humano. El sistema contempla tres canales de entrada (correo electronico, formulario web y llamada telefonica con transcripcion automatica) y tres sectores de derivacion (Sistemas, Operaciones y Soporte Tecnico). La validacion experimental se realiza sobre un corpus en idioma espanol rioplatense recolectado en una unica organizacion del sector servicios.

Quedan fuera del alcance los siguientes aspectos, los cuales se proponen como lineas de trabajo futuro en el Capitulo~10. Primero, la integracion con sistemas externos de inventario de activos informaticos o con sistemas de gestion de configuracion. Segundo, los modelos de aprendizaje supervisado entrenados con corpus propios de la organizacion, dado que el volumen actual de datos etiquetados resulta insuficiente para esa estrategia. Tercero, la implementacion de paneles de monitoreo avanzados con visualizacion en tiempo real de metricas operativas. Cuarto, la automatizacion de la etapa de resolucion tecnica, que trasciende el alcance del presente trabajo.

Las delimitaciones del estudio incluyen: la restriccion geografica a una unica organizacion de la provincia de Mendoza, lo que limita la validez externa de los hallazgos; la dependencia operativa del servicio externo de inferencia Gemini~2.5 Flash, cuyo rendimiento y condiciones comerciales escapan al control de los investigadores; y la naturaleza temporal de los resultados, que reflejan el comportamiento del modelo en la version vigente al momento del estudio (junio de 2026), sujeta a modificaciones futuras por parte del proveedor.

\newpage
